\documentclass[addpoints,spanish, 12pt,a4paper]{exam}
%\documentclass[answers, spanish, 12pt,a4paper]{exam}
\printanswers
\pointpoints{punto}{puntos}
\hpword{Puntos:}
\vpword{Puntos:}
\htword{Total}
\vtword{Total}
\hsword{Resultado:}
\hqword{Ejercicio:}
\vqword{Ejercicio:}

\usepackage[utf8]{inputenc}
\usepackage[spanish]{babel}
\usepackage{eurosym}
%\usepackage[spanish,es-lcroman, es-tabla, es-noshorthands]{babel}


\usepackage[margin=1in]{geometry}
\usepackage{amsmath,amssymb}
\usepackage{multicol}
\usepackage{yhmath}



\usepackage{graphicx}
\graphicspath{{../img/}} 

\newcommand{\class}{4º Académicas}
\newcommand{\examdate}{\today}
\newcommand{\examnum}{Examen de potencias, radicales y logaritmos}
\newcommand{\tipo}{A}


\newcommand{\timelimit}{50 minutos}

\renewcommand{\solutiontitle}{\noindent\textbf{Solución:}\enspace}


\pagestyle{head}
\firstpageheader{\includegraphics[width=0.2\columnwidth]{header_left}}{\textbf{Departamento de Matemáticas\linebreak \class}\linebreak \examnum}{\includegraphics[width=0.1\columnwidth]{header_right}}
\runningheader{\class}{\examnum}{Página \thepage\ of \numpages}
\runningheadrule


\begin{document}

\noindent
\begin{tabular*}{\textwidth}{l @{\extracolsep{\fill}} r @{\extracolsep{6pt}} }
\textbf{Nombre:} \makebox[3.5in]{\hrulefill} & \textbf{Fecha:}\makebox[1in]{\hrulefill} \\
 & \\
\textbf{Tiempo: \timelimit} & Tipo: \tipo 
\end{tabular*}
\rule[2ex]{\textwidth}{2pt}
Esta prueba tiene \numquestions\ ejercicios. La puntuación máxima es de \numpoints. 
La nota final de la prueba será la parte proporcional de la puntuación obtenida sobre la puntuación máxima. %Para la recuperación de pendientes de 3º se tendrán en cuenta los apartados: 3.a, 3.b, 4, 5.a, 5.b, 5.c

\begin{center}


\addpoints
 %\gradetable[h][questions]
	\pointtable[h][questions]
\end{center}

\noindent
\rule[2ex]{\textwidth}{2pt}

\begin{questions}

\question[2] Indica a cuáles de los conjuntos
$\mathbb{N}$, $\mathbb{Z}$, $\mathbb{Q}$, $\mathbb{R}$ pertenecen cada uno de los siguientes números:
\begin{center}
\begin{tabular}{|c |c |c |c |c|}\hline
    &$\mathbb{N}$& $\mathbb{Z}$& $\mathbb{Q}$&$\mathbb{R}$\\ 
    \hline $(-1)^3$&&&&\\ \hline $-\frac{20}{2}$&&&&\\ \hline $-4^2$ &&&&\\ \hline
    $1,\wideparen{3}$&&&&\\
\hline
$-\frac{16}{4}$&&&&\\
\hline
$-\sqrt{25}$&&&&\\
\hline
$\sqrt{8}$&&&&\\
\hline
$4$&&&&\\
\hline
$\pi$&&&&\\
\hline
    \end{tabular}
\end{center}

\begin{solution}
\begin{tabular}{|c |c |c |c |c|}\hline
    &$\mathbb{N}$& $\mathbb{Z}$& $\mathbb{Q}$&$\mathbb{R}$\\ 
    \hline$(-1)^3$ & False & True & True & True\\ \hline $-\frac{20}{2}$ & False & True & True & True\\ \hline $-4^2$ & False & True & True & True\\ \hline \end{tabular}
\end{solution}

\addpoints

% \question[1] Representa en la recta real y en forma de intervalo el siguiente conjunto numérico:
% \addpoints % to omit double points count
% $$\left\{ x \in \mathbb{R} \left| \ 8 \leqslant x  \right. \right\}$$

% \begin{solution}
% $$\left[8 \ , +\infty\right)$$ 
% \end{solution}

\question[1] Calcula $A\cup B$ y $A\cap B$ si  $A=\left[-1, 4\right)$ y $B=\{x|1 < x \leq 4\}$ (da el resultado en forma de intervalo y de desigualdad):
\addpoints % to omit double points count

\begin{solution}
Unión: \\ $A\cup B=\left[-1, 4\right]$ ó $A\cup B= \{x|-1 \leq x \wedge x \leq 4\}$ \\Intersección: \\ $A\cap B =\left(1, 4\right)$ ó $A\cap B=\{x|1 < x \wedge x < 4\}$
\end{solution}


\question Opera:
%\noaddpoints % to omit double points count

\begin{parts}
\part[1] \[\dfrac{\left(2^3\cdot3^2\cdot5\right)^{-4}}{\left(2^{-2}\cdot3^{-3}\right)^{3}}\]
\begin{solution}
$=\dfrac{2^{-12}\cdot3^{-8}\cdot5^{-4}}{2^{-6}\cdot3^{-9}}=\dfrac{3}{2^6\cdot5^4}$ 
\end{solution}
% \part[1] \[\left(\dfrac{6p^3d^2}{5q}\right)^4\cdot\left(\dfrac{20p^2q^3}{24d}\right)^4\]
% \begin{solution}
% $=\left(\dfrac{6p^3d^2}{5q}\cdot\dfrac{20p^2q^3}{24d}\right)^4=\left(p^5dq^2\right)^4=p^{20}d^4q^8 $
% \end{solution}
\end{parts}

\addpoints

\question[2] Expresa en notación científica, opera y simplifica:
\addpoints 

$$\dfrac{0'0012\cdot0'002\cdot100000}{8000\cdot0'0003\cdot0'01}$$
\begin{solution} $=\dfrac{1,2\cdot10^{-3}\cdot2\cdot10^{-3}\cdot10^{5}}{8\cdot10^{3}\cdot3\cdot10^{-4}\cdot10^{-2}}=\dfrac{2,4\cdot10^{-1}}{24\cdot10^{-3}}=10$ \end{solution}
 

\question Opera y simplifica cada una de estas expresiones:
%\noaddpoints % to omit double points count

\begin{parts}
\part[1] \[4\sqrt{20}-3\sqrt{45}+11\sqrt{125}-20\sqrt{5}\]
\begin{solution}
$=4\cdot2\sqrt{5}-3\cdot3\sqrt{5}+11\cdot5\sqrt{5}-20\sqrt{5}=\left(8-9+55-20\right)\sqrt{5}=34\sqrt{5}$
\end{solution}

\part[1] \[\sqrt{72}\cdot3\sqrt{8}\]
\begin{solution}
$=\sqrt{72}\cdot\sqrt{72}=\sqrt{72^2}=72$
\end{solution}

\part[1] \[\sqrt[4]{\frac{25}{9}\sqrt[3]{\frac{9}{25}}}\]
\begin{solution}
$=\sqrt[4]{\sqrt[3]{\left(\dfrac{25}{9}\right)^3\dfrac{9}{25}}}=\sqrt[12]{\left(\dfrac{25}{9}\right)^2}=\sqrt[6]{\left(\dfrac{5}{3}\right)^2}=\sqrt[3]{\frac{5}{3}} $
\end{solution}

\part[2] \[\dfrac{\left(3\sqrt{2}+\sqrt{3}\right)^2}{3}\]
\begin{solution}
$ =\dfrac{9\cdot2+3+2\cdot3\sqrt{2}\sqrt{3}}{3}=3\cdot2+1+2\sqrt{2}\sqrt{3}=7+2\sqrt{6}$
\end{solution}


\end{parts}
\addpoints

\question Racionaliza y simplifica:
%\noaddpoints % to omit double points count
\begin{parts}
\part[2] \[\dfrac{10}{2\sqrt{3}-\sqrt{2}}\]
\begin{solution}
$=\dfrac{10\cdot\left(2\sqrt{3}+\sqrt{2}\right)}{\left(2\sqrt{3}-\sqrt{2}\right)\cdot\left(2\sqrt{3}+\sqrt{2}\right)}=\dfrac{10\cdot\left(2\sqrt{3}+\sqrt{2}\right)}{4\cdot3-2}=2\sqrt{3}+\sqrt{2}$
\end{solution}


\part[2] \[\dfrac{\sqrt{45}+\sqrt{180}}{\sqrt{176}+4\sqrt{44}}\]
\begin{solution}
$=\dfrac{3\sqrt{5}+3\cdot2\sqrt{5}}{4\sqrt{11}+4\cdot2\sqrt{11}}=\dfrac{9\sqrt{5}}{12\sqrt{11}}=\dfrac{3\sqrt{55}}{44}$
\end{solution}



\end{parts}
\addpoints


\question Sin utilizar la calculadora, resuelve los siguientes logaritmos:
\begin{multicols}{2}

\begin{parts}

\part[1] ${\log _3}27$ \begin{solution} ${\log _3}27 = {\log _3}{3^3} = 3$ \end{solution}

\part[1] ${\log _{{\raise0.7ex\hbox{$1$} \!\mathord{\left/
 {\vphantom {1 3}}\right.\kern-\nulldelimiterspace}
\!\lower0.7ex\hbox{$3$}}}}\frac{1}{{81}}$ \begin{solution} ${\log _{{\raise0.7ex\hbox{$1$} \!\mathord{\left/
 {\vphantom {1 3}}\right.\kern-\nulldelimiterspace}
\!\lower0.7ex\hbox{$3$}}}}\frac{1}{{81}} = {\log _{{\raise0.7ex\hbox{$1$} \!\mathord{\left/
 {\vphantom {1 3}}\right.\kern-\nulldelimiterspace}
\!\lower0.7ex\hbox{$3$}}}}{\left( {\frac{1}{3}} \right)^4} = 4$\end{solution}

\part[1] ${\log _3}\frac{1}{{81}}$ \begin{solution} ${\log _3}\frac{1}{{81}} = {\log _3}{3^{ - 4}} =  - 4$ \end{solution}

\part[1] ${\log _5}\sqrt {125} $ \begin{solution} ${\log _5}\sqrt {125}  = {\log _5}{5^{{\raise0.7ex\hbox{$3$} \!\mathord{\left/
 {\vphantom {3 2}}\right.\kern-\nulldelimiterspace}
\!\lower0.7ex\hbox{$2$}}}} = \frac{3}{2}$ \end{solution}
\end{parts}
\end{multicols}


\addpoints


\question Calcula:
%\noaddpoints % to omit double points count

\begin{parts}


% \part[1] \[\log_3 \frac{1}{9} -  \log_5 0,2 +\log_6 \frac{1}{36} - \log_2 0,5\]
% \begin{solution}
% $=-2-\left(-1\right)+\left(-2\right)-\left(-1\right)=-2$
% \end{solution}

\part[1] $\log_3 \frac{1}{27} -  \log_5 0.2 +\log_7 \frac{1}{49} - \log_2 0.5$
\begin{solution}
$\log_3 \frac{1}{27} -  \log_5 0.2 +\log_7 \frac{1}{49} - \log_2 0.5= (-3.0 )+( 1.0 )+( -2.0 )+( 1.0)=-3.0$
\end{solution}




\end{parts}

\addpoints


\question Calcula aplicando la propiedades de los logaritmos:
%\noaddpoints % to omit double points count
$\log_2 20 + 2\log_2 3 - \log_2 5 - \log_2 9$
$2$
\begin{parts}
\part[2] $\log_2 20 + 2\log_2 3 - \log_2 5 - \log_2 9$
\begin{solution}
$=\log_2 \left(\dfrac{20\cdot3^2}{5\cdot9}\right)=\log_2 4=2$
\end{solution}

\end{parts}

\addpoints

\end{questions}

\end{document}
\grid
